\chapter{De los datos al modelo y del modelo a los datos}\label{Sim_diarias}
\chapterquote{}{}
\section{Introducción}\label{P:introduccion}

A lo largo de la historia, el ser humano ha intentado comprender el mundo que lo rodea. Para ello existen dos metodologías predominantes en el ámbito, la rama 
cualitativa, predominante en la biología y la rama cuantitativa, predominante en la física. Ante esto, no es de extrañar la aparición de gente que busque 
unificar estos rubros usando metodologías aprendidas en una de las ramas, aplicándolas sobre objetos de estudios de la otra. En lo que a nuestro trabajo 
concierne, esta aproximación de ambos mundos toma el nombre de biología o comportamiento cuantitativo. Parte de lo que se estudia en este mundo interdisciplinario
son los modelos de caminantes, los cuales conectan conceptos vistos comúnmente en materias de probabilidad y procesos aleatorios dentro de la carrera de física, que 
podemos usar para intentar replicar los patrones de movimiento de especies registrados con localizadores de GPS.

Los procesos de caminantes aleatorios llevan tiempo siendo estudiados, desde los caminantes unidimensionales hasta caminantes bidimensionales con 
estados internos que definen el comportamiento de los individuos \cite{codling2008}. Este tipo de modelos ofrece un marco potente para describir una 
amplia variedad de fenómenos biológicos, desde la dinámica poblacional hasta el movimiento celular, y han sido ampliamente utilizados para analizar 
trayectorias de animales obtenidas mediante dispositivos de GPS \cite{bartumeus2005}.

Dentro de este último subconjunto de modelos, vamos a tomar como punto de partida el modelo de \textit{Stop and Go} \cite{peruani2023stopandgo}, en donde 
los individuos alternan entre dos estados posibles: uno asociado al movimiento y otro al reposo. La principal diferencia de este modelo respecto a otros 
populares, como el \textit{run-and-tumble} \cite{codling2008}, radica en el rol que cumple el estado de reposo: mientras que en el run-and-tumble la 
reorientación del individuo ocurre durante esta fase, en el Stop and Go el individuo, al encontrarse en reposo, deja pasar el tiempo sin realizar ninguna 
acción más allá de definir si permanece en ese estado o cambia al de movimiento.

Nuestro objetivo en este capítulo es construir un modelo de tipo \textit{Stop and Go} capaz de reproducir el comportamiento diario de las tortugas, calibrando 
sus parámetros a partir de los datos de GPS descritos en el Capítulo \ref{Introduccion}. En particular, buscamos que las trayectorias simuladas reproduzcan 
las distribuciones de tiempos de permanencia en cada estado (movimiento y reposo) reportadas por C. Callenbach \cite{callenbach2025}, así como otras métricas
estadísticas relevantes, como el desplazamiento cuadrático medio (MSD) y la distribución de desplazamientos entre dos adquisiciones de ubicación del dispositivo de GPS.


\section{Nuestro modelo}\label{P:Mi_Modelo}

El modelo desarrollado en este capítulo describe el movimiento de cada tortuga como un caminante aleatorio de tipo \textit{Stop and Go} que se desplaza en 
un espacio homogéneo continuo, es decir, sin refugios ni otras estructuras que introduzcan heterogeneidad espacial. Simulamos $N \approx 100.000$ individuos 
de forma independiente, sin interacción alguna entre ellos, cada uno alternando entre dos estados: movimiento (\textit{Go}) y reposo (\textit{Stop}). 
El ciclo día-noche se incorpora únicamente como una ventana horaria activa, la cual abarca desde las 8hs hasta las 20hs, fuera de este intervalo, 
el individuo permanece inmóvil. En estas simulaciones tomamos los pasos temporales equivalentes a $1.7 s$ (tiempo entre mediciones de acelerómetro) y 
se dejó evolucionar hasta llegar a un equivalente de 2 meses. La selección de la duración de la simulación fue inspirada en el comportamiento no brumativo de 
las tortugas, que se detallará más adelante, en el Cápitulo \ref{Sim_año}, cuando se analice el comportamiento de esta especie a lo largo del año.


Para determinar cuánto tiempo permanece un individuo en cada uno de los dos estados, existen al menos dos estrategias posibles. La primera consiste en 
definir una tasa de transición instantánea $\gamma(t)$, función del tiempo ya transcurrido en el estado actual, y evaluarla en cada paso de tiempo discreto 
$\delta t$ de la simulación para decidir si el individuo cambia de estado. Si bien este enfoque permite incorporar naturalmente una tasa que varíe con el 
tiempo, lo cual incrementa considerablemente el costo computacional al simular una gran cantidad de individuos.

La segunda estrategia, que es la que adoptamos en esta tesis, consiste en sortear directamente, al momento de ingresar a un estado, el tiempo total de 
permanencia en él, utilizando la distribución empírica completa reportada por Callenbach \cite{callenbach2025}. Para esto empleamos el método de la 
transformada inversa (\textit{inverse-CDF sampling}): dada una densidad de probabilidad $f(t)\sim t^{-\alpha}$ truncada entre un tiempo mínimo $t_{min}$ 
y un tiempo máximo $t_{max}$, es posible invertir su función de distribución acumulada y obtener una expresión cerrada que, a partir de un número aleatorio 
$u$ uniformemente distribuido en $[0,1]$, entrega directamente el tiempo sorteado:

\begin{equation}
    t = \left[(t_{max}^{\,1-\alpha} - t_{min}^{\,1-\alpha})\,u + t_{min}^{\,1-\alpha}\right]^{\frac{1}{1-\alpha}}
\end{equation}

Este método permite obtener, en una única operación, el tiempo completo que un individuo permanecerá en el estado de movimiento o de reposo, sin necesidad 
de evaluar la probabilidad de transición paso a paso. Dado que trabajamos con simulaciones de una gran cantidad de individuos y de larga duración temporal, 
esta reducción en el costo computacional resulta determinante para la viabilidad del trabajo.


En el caso de querer realizar cálculos teóricos, puede ser de utilidad tener una ecuación para la tasa de transición $\gamma$ de un estado a otro, en función del 
tiempo transcurrido en el estado actual, $t$, teniendo en consideración los pasos temporales discretos de la simulación, $\delta t$ y una distribución 
temporal objetivo tipo ley de potencia $\propto t^{-\alpha}$. El desarrollo paraobtener esta expresión se puede encontrar en el apéndice \ref{A:transition rate PL}, 
dando como resultado la ecuación:

\begin{equation}
    \gamma (t)=\frac{(t+\delta t)^{1-\alpha}-t^{1-\alpha}}{t_{\text{max}}^{1-\alpha}-t^{1-\alpha}} 
\end{equation}

\subsection*{Stop and Go}

Existen distintos métodos de implementación del modelo de caminante \textit{Stop and Go}. Ya vimos algunas variaciones respecto a los ratio de transición.
Sin embargo, también se pueden variar cosas en funcionamiento del estado de \textit{Go}. El modelo utilizado se conoce en la bibliografía como 
\textit{Correlated Random Walk}, en donde existe una correlación entre dos pasos consecutivos. En este caso, la correlación aparece al variar la dirección
en cada paso temporal, donde el nuevo ángulo de orientación $\theta_{i+1}$ se obtiene a partir del ángulo del paso anterior $\theta_i$, 
sumándole un ruido gaussiano de media nula y desviación estándar $\sigma_\theta$,

\begin{equation}
    \theta_{i+1} = \theta_i + \eta_i, \qquad \eta_i \sim \mathcal{N}(0,\sigma_\theta^2).
\end{equation}

De esta manera, una vez tenida la nueva dirección, el individuo se desplaza en $l = 1 u.a.$ y la nueva posición del individuo queda definida como:

\begin{equation}
    x_{i+1} = x_i + l\cos(\theta_{i+1}), \qquad y_{i+1} = y_i + l\sin(\theta_{i+1}).
\end{equation}

Esta correlación angular introduce persistencia direccional en la trayectoria: a diferencia de un caminante aleatorio simple, en donde la dirección se 
sortea de forma independiente en cada paso, el individuo tiende a mantener su rumbo a corto plazo, generando trayectorias localmente rectas que se 
curvan gradualmente a medida que se acumulan las perturbaciones angulares. Estas persistencias son más fuertez cuanto menor sea el valor de $\sigma_\theta$, 
el cual es un parámetro del sistema. 

En el estado de \textit{Stop}, por el contrario, la posición del individuo permanece fija: $x_{i+1}=x_i$, $y_{i+1}=y_i$, y el único proceso que avanza es 
el conteo del tiempo.

La duración de cada \textit{bout} (período temporal ininterrumpido que pasa en un estado), tanto de \textit{Go} como de \textit{Stop}, se sortea al momento 
de ingresar al estado correspondiente. En el análisis de este trabajo se utilizaron tanto duraciones constantes en cada uno de los estados, como las distribuciones 
tipo ley de potencias, inspiradas en las distribuciones empíricas reportadas por Callenbach \cite{callenbach2025} para cada uno de los dos estados.


\section{Resultados del modelo}\label{P:Resultados_diario}













\section{Discusiones y conlcusiones}\label{P:Conclusiones_diario}




%%% Local Variables: 
%%% mode: latex
%%% TeX-master: "template"
%%% End: 
